\title{2024 年高考数学 全国甲卷-理科}
\date{}
\maketitle

\section*{一、选择题}
本题共12小题，每小题5分，共60分，在每小题给出的四个选项中，只有一项是符合题目要求的。

\begin{question}[points=5]
设 $z=5+\mathrm{i}$，则 $\mathrm{i}(z+\overline{z})=$ （   ）
\begin{choices}
  \item $10\mathrm{i}$
  \item $2\mathrm{i}$
  \item $10$
  \item $-2$
\end{choices}
\begin{solution}
\textbf{答案：}A

\textbf{分析：}结合共轭复数与复数的基本运算直接求解。

\textbf{详解：}由 $z=5+\mathrm{i}\Rightarrow\overline{z}=5-\mathrm{i},z+\overline{z}=10$，则 $\mathrm{i}(z+\overline{z})=10\mathrm{i}$。
\end{solution}
\end{question}

\begin{question}[points=5]
集合 $A=\{1,2,3,4,5,9\},B=\{x|x\in A\}$，则 $\complement_A(A\cap B)=$ （   ）
\begin{choices}
  \item $\{1,4,9\}$
  \item $\{3,4,9\}$
  \item $\{1,2,3\}$
  \item $\{2,3,5\}$
\end{choices}
\begin{solution}
\textbf{答案：}D

\textbf{分析：}由集合 $B$ 的定义求出 $B$，结合交集与补集运算即可求解。

\textbf{详解：}因为 $A=\{1,2,3,4,5,9\},B=\{x|x\in A\}$，所以 $B=\{1,4,9,16,25,81\}$，则 $A\cap B=\{1,4,9\}$，$\complement_A(A\cap B)=\{2,3,5\}$。
\end{solution}
\end{question}

\begin{question}[points=5]
若实数 $x,y$ 满足约束条件 $\begin{cases} 4x-3y-3\ge0 \\ x-2y-2\le0 \\ 2x+6y-9\le0 \end{cases}$，则 $z=x-5y$ 的最小值为（   ）

\begin{center}
  \includegraphics[width=.38\linewidth]{question-3-1.jpg}
\end{center}
\begin{choices}
  \item $5$
  \item $\dfrac{1}{2}$
  \item $-2$
  \item $-\dfrac{7}{2}$
\end{choices}
\begin{solution}
\textbf{答案：}D

\textbf{分析：}画出可行域后，利用 $z$ 的几何意义计算即可得。

\textbf{详解：}作出可行域如图，由 $z=x-5y$ 可得 $y=\frac{1}{5}x-\frac{z}{5}$，则 $z$ 的几何意义为直线 $y=\frac{1}{5}x-\frac{z}{5}$ 的截距的 $-\frac{1}{5}$，则该直线截距取最大值时，$z$ 有最小值，此时直线过点 $A$，联立 $\begin{cases}4x-3y-3=0\\2x+6y-9=0\end{cases}$，解得 $A\left(\frac{3}{2},1\right)$，则 $z_{\min}=\frac{3}{2}-5\times1=-\frac{7}{2}$。
\end{solution}
\end{question}

\begin{question}[points=5]
等差数列 $\{a_n\}$ 的前 $n$ 项和为 $S_n$，若 $S_5=S_{10}$，$a_5=1$，则 $a_1=$ （   ）
\begin{choices}
  \item $-2$
  \item $\dfrac{7}{3}$
  \item $1$
  \item $2$
\end{choices}
\begin{solution}
\textbf{答案：}B

\textbf{分析：}由 $S_5=S_{10}$ 结合等差中项的性质可得 $a_8=0$，即可计算出公差，即可得 $a_1$ 的值。

\textbf{详解：}由 $S_{10}-S_5=a_6+a_7+a_8+a_9+a_{10}=5a_8=0$，则 $a_8=0$，则等差数列 $\{a_n\}$ 的公差 $d=\frac{a_8-a_5}{3}=-\frac{1}{3}$，故 $a_1=a_5-4d=1-4\times\left(-\frac{1}{3}\right)=\frac{7}{3}$。
\end{solution}
\end{question}

\begin{question}[points=5]
已知双曲线的两个焦点分别为 $(0,4),(0,-4)$，点 $(-6,4)$ 在该双曲线上，则该双曲线的离心率为（   ）
\begin{choices}
  \item $4$
  \item $3$
  \item $2$
  \item $\sqrt{2}$
\end{choices}
\begin{solution}
\textbf{答案：}C

\textbf{分析：}由焦点坐标可得焦距 $2c$，结合双曲线定义计算可得 $2a$，即可得离心率。

\textbf{详解：}设 $F_1(0,-4),F_2(0,4),P(-6,4)$，则 $|F_1F_2|=2c=8$，$|PF_1|=\sqrt{6^2+(4+4)^2}=10$，$|PF_2|=\sqrt{6^2+(4-4)^2}=6$，则 $2a=|PF_1|-|PF_2|=10-6=4$，则 $e=\frac{2c}{2a}=\frac{8}{4}=2$。
\end{solution}
\end{question}

\begin{question}[points=5]
设函数 $f(x)=\dfrac{e^x+2\sin x}{1+x^2}$，则曲线 $y=f(x)$ 在 $(0,1)$ 处的切线与两坐标轴围成的三角形的面积为（   ）
\begin{choices}
  \item $\dfrac{1}{6}$
  \item $\dfrac{1}{3}$
  \item $\dfrac{1}{2}$
  \item $\dfrac{2}{3}$
\end{choices}
\begin{solution}
\textbf{答案：}A

\textbf{分析：}借助导数的几何意义计算可得其在点 $(0,1)$ 处的切线方程，即可得其与坐标轴交点坐标，即可得其面积。

\textbf{详解：}$f'(x)=\frac{(e^x+2\cos x)(1+x^2)-(e^x+2\sin x)\cdot2x}{(1+x^2)^2}$，则 $f'(0)=\frac{(e^0+2\cos0)(1+0)-(e^0+2\sin0)\times0}{(1+0)^2}=3$，即该切线方程为 $y-1=3x$，即 $y=3x+1$，令 $x=0$，则 $y=1$，令 $y=0$，则 $x=-\frac{1}{3}$，故该切线与两坐标轴所围成的三角形面积 $S=\frac{1}{2}\times1\times\left|-\frac{1}{3}\right|=\frac{1}{6}$。
\end{solution}
\end{question}

\begin{question}[points=5]
函数 $f(x)=-x^2+(e^x-e^{-x})\sin x$ 在区间 $[-2.8,2.8]$ 的大致图像为（   ）

\begin{center}
  \includegraphics[width=.38\linewidth]{question-7-1.jpg}
\end{center}
\begin{choices}
  \item A
  \item B
  \item C
  \item D
\end{choices}
\begin{solution}
\textbf{答案：}B

\textbf{分析：}利用函数的奇偶性可排除 A、C，代入 $x=1$ 可得 $f(1)>0$，可排除 D。

\textbf{详解：}$f(-x)=-(-x)^2+(e^{-x}-e^x)\sin(-x)=-x^2+(e^x-e^{-x})\sin x=f(x)$，又函数定义域为 $[-2.8,2.8]$，故该函数为偶函数，可排除 A、C，又 $f(1)=-1+\left(e-\frac{1}{e}\right)\sin1>-1+\left(e-\frac{1}{e}\right)\sin\frac{\pi}{6}=\frac{e}{2}-1-\frac{1}{2e}>\frac{1}{4}-\frac{1}{2e}>0$，故可排除 D。
\end{solution}
\end{question}

\begin{question}[points=5]
已知 $\dfrac{\cos\alpha}{\cos\alpha-\sin\alpha}=3$，则 $\tan\left(\alpha+\dfrac{\pi}{4}\right)=$ （   ）
\begin{choices}
  \item $2\sqrt{3}+1$
  \item $2\sqrt{3}-1$
  \item $\dfrac{3}{2}$
  \item $1-\sqrt{3}$
\end{choices}
\begin{solution}
\textbf{答案：}B

\textbf{分析：}先将 $\dfrac{\cos\alpha}{\cos\alpha-\sin\alpha}$ 弦化切求得 $\tan\alpha$，再根据两角和的正切公式即可求解。

\textbf{详解：}因为 $\dfrac{\cos\alpha}{\cos\alpha-\sin\alpha}=3$，所以 $\dfrac{1}{1-\tan\alpha}=3$，$\Rightarrow\tan\alpha=1-\frac{1}{3}$，所以 $\tan\left(\alpha+\frac{\pi}{4}\right)=\frac{\tan\alpha+1}{1-\tan\alpha}=2\sqrt{3}-1$。
\end{solution}
\end{question}

\begin{question}[points=5]
已知向量 $\vec{a}=(x+1,x),\vec{b}=(x,2)$，则（   ）
\begin{choices}
  \item “$x=-3$”是“$\vec{a}\perp\vec{b}$”的必要条件
  \item “$x=-3$”是“$\vec{a}\parallel\vec{b}$”的必要条件
  \item “$x=0$”是“$\vec{a}\perp\vec{b}$”的充分条件
  \item “$x=-1+\sqrt{3}$”是“$\vec{a}\parallel\vec{b}$”的充分条件
\end{choices}
\begin{solution}
\textbf{答案：}C

\textbf{分析：}根据向量垂直和平行的坐标表示即可得到方程，解出即可。

\textbf{详解：}对 A，当 $\vec{a}\perp\vec{b}$ 时，则 $\vec{a}\cdot\vec{b}=0$，所以 $x\cdot(x+1)+2x=0$，解得 $x=0$ 或 $-3$，即必要性不成立，故 A 错误；对 C，当 $x=0$ 时，$\vec{a}=(1,0),\vec{b}=(0,2)$，故 $\vec{a}\cdot\vec{b}=0$，所以 $\vec{a}\perp\vec{b}$，即充分性成立，故 C 正确；对 B，当 $\vec{a}\parallel\vec{b}$ 时，则 $2(x+1)=x^2$，解得 $x=1\pm\sqrt{3}$，即必要性不成立，故 B 错误；对 D，当 $x=-1+\sqrt{3}$ 时，不满足 $2(x+1)=x^2$，所以 $\vec{a}\parallel\vec{b}$ 不成立，即充分性不立，故 D 错误。
\end{solution}
\end{question}

\begin{question}[points=5]
设 $\alpha,\beta$ 是两个平面，$m,n$ 是两条直线，且 $\alpha\cap\beta=m$。下列四个命题：

①若 $m\parallel n$，则 $n\parallel\alpha$ 或 $n\parallel\beta$

②若 $m\perp n$，则 $n\perp\alpha$，$n\perp\beta$

③若 $n\parallel\alpha$，且 $n\parallel\beta$，则 $m\parallel n$

④若 $n$ 与 $\alpha$ 和 $\beta$ 所成的角相等，则 $m\perp n$

其中所有真命题的编号是（   ）
\begin{choices}
  \item ①③
  \item ②④
  \item ①②③
  \item ①③④
\end{choices}
\begin{solution}
\textbf{答案：}A

\textbf{分析：}根据线面平行的判定定理即可判断①；举反例即可判断②④；根据线面平行的性质即可判断③。

\textbf{详解：}对①，当 $n\subset\alpha$，因为 $m\parallel n$，$m\subset\beta$，则 $n\parallel\beta$，当 $n\subset\beta$，因为 $m\parallel n$，$m\subset\alpha$，则 $n\parallel\alpha$，当 $n$ 既不在 $\alpha$ 也不在 $\beta$ 内，因为 $m\parallel n$，$m\subset\alpha$，$m\subset\beta$，则 $n\parallel\alpha$ 且 $n\parallel\beta$，故①正确；对②，若 $m\perp n$，则 $n$ 与 $\alpha,\beta$ 不一定垂直，故②错误；对③，过直线 $n$ 分别作两平面与 $\alpha,\beta$ 分别相交于直线 $s$ 和直线 $t$，因为 $n\parallel\alpha$，过直线 $n$ 的平面与平面 $\alpha$ 的交线为直线 $s$，则根据线面平行的性质定理知 $n\parallel s$，同理可得 $n\parallel t$，则 $s\parallel t$，因为 $s\not\subset$ 平面 $\beta$，$t\subset$ 平面 $\beta$，则 $s\parallel$ 平面 $\beta$，因为 $s\subset$ 平面 $\alpha$，$\alpha\cap\beta=m$，则 $s\parallel m$，又因为 $n\parallel s$，则 $m\parallel n$，故③正确；对④，若 $\alpha\cap\beta=m$，$n$ 与 $\alpha$ 和 $\beta$ 所成的角相等，如果 $n\parallel\alpha$，$n\parallel\beta$，则 $m\parallel n$，故④错误；综上只有①③正确。
\end{solution}
\end{question}

\begin{question}[points=5]
在 $\triangle ABC$ 中内角 $A,B,C$ 所对边分别为 $a,b,c$，若 $B=\dfrac{\pi}{3}$，$b^2=\dfrac{9}{4}ac$，则 $\sin A+\sin C=$ （   ）
\begin{choices}
  \item $\dfrac{3}{2}$
  \item $\sqrt{2}$
  \item $\dfrac{\sqrt{7}}{2}$
  \item $\dfrac{3}{\sqrt{2}}$
\end{choices}
\begin{solution}
\textbf{答案：}C

\textbf{分析：}利用正弦定理得 $\sin A\sin C=\frac{1}{3}$，再利用余弦定理有 $a^2+c^2=\frac{13}{4}ac$，再利用正弦定理得到 $\sin^2 A+\sin^2 C$ 的值，最后代入计算即可。

\textbf{详解：}因为 $B=\frac{\pi}{3},b^2=\frac{9}{4}ac$，则由正弦定理得 $\sin A\sin C=\frac{4}{9}\sin^2 B=\frac{1}{3}$。由余弦定理可得 $b^2=a^2+c^2-ac=\frac{9}{4}ac$，即 $a^2+c^2=\frac{13}{4}ac$，根据正弦定理得 $\sin^2 A+\sin^2 C=\frac{13}{4}\sin A\sin C=\frac{13}{12}$，所以 $(\sin A+\sin C)^2=\sin^2 A+\sin^2 C+2\sin A\sin C=\frac{7}{4}$，因为 $A,C$ 为三角形内角，则 $\sin A+\sin C>0$，则 $\sin A+\sin C=\frac{\sqrt{7}}{2}$。
\end{solution}
\end{question}

\begin{question}[points=5]
已知 $b$ 是 $a,c$ 的等差中项，直线 $ax+by+c=0$ 与圆 $x^2+y^2+4y-1=0$ 交于 $A,B$ 两点，则 $|AB|$ 的最小值为（   ）
\begin{choices}
  \item $2$
  \item $3$
  \item $4$
  \item $2\sqrt{5}$
\end{choices}
\begin{solution}
\textbf{答案：}C

\textbf{分析：}结合等差数列性质将 $c$ 代换，求出直线恒过的定点，采用数形结合法即可求解。

\textbf{详解：}因为 $a,b,c$ 成等差数列，所以 $2b=a+c$，$c=2b-a$，代入直线方程 $ax+by+c=0$ 得 $ax+by+2b-a=0$，即 $a(x-1)+b(y+2)=0$，令 $\begin{cases}x-1=0\\y+2=0\end{cases}$ 得 $\begin{cases}x=1\\y=-2\end{cases}$，故直线恒过 $(1,-2)$，设 $P(1,-2)$，圆化为标准方程得 $C:x^2+(y+2)^2=5$，设圆心为 $C$，画出直线与圆的图形，由图可知，当 $PC\perp AB$ 时，$|AB|$ 最小，$|PC|=1,|AC|=r=\sqrt{5}$，此时 $|AB|=2|AP|=2\sqrt{|AC|^2-|PC|^2}=2\sqrt{5-1}=4$。
\end{solution}
\end{question}

\section*{二、填空题}
本题共4小题，每小题5分，共20分。

\begin{question}[points=5]
$\left(\dfrac{1}{3}+x\right)^{10}$ 的展开式中，各项系数的最大值是。
\fillin[{$5$}]
\begin{solution}
\textbf{答案：}$5$

\textbf{分析：}先设展开式中第 $r+1$ 项系数最大，则根据通项公式有不等式组，进而求出 $r$ 即可求解。

\textbf{详解：}由题展开式通项公式为 $T_{r+1}=C_{10}^r\left(\frac{1}{3}\right)^{10-r}x^r$，$0\le r\le10$ 且 $r\in\mathbb{Z}$，设展开式中第 $r+1$ 项系数最大，则 $\begin{cases} C_{10}^r\left(\frac{1}{3}\right)^{10-r}\ge C_{10}^{r+1}\left(\frac{1}{3}\right)^{9-r} \\ C_{10}^r\left(\frac{1}{3}\right)^{10-r}\ge C_{10}^{r-1}\left(\frac{1}{3}\right)^{11-r} \end{cases}$，$\Rightarrow\begin{cases} r\ge\frac{29}{4} \\ r\le\frac{33}{4} \end{cases}$，又 $r\in\mathbb{Z}$，故 $r=8$，所以展开式中系数最大的项是第9项，且该项系数为 $C_{10}^8\left(\frac{1}{3}\right)^2=5$。
\end{solution}
\end{question}

\begin{question}[points=5]
已知甲、乙两个圆台上、下底面的半径均为 $r_1$ 和 $r_2$，母线长分别为 $2(r_2-r_1)$ 和 $3(r_2-r_1)$，则两个圆台的体积之比 $\dfrac{V_{\text{甲}}}{V_{\text{乙}}}=$ 。
\fillin[{$\dfrac{\sqrt{6}}{4}$}]
\begin{solution}
\textbf{答案：}$\dfrac{\sqrt{6}}{4}$

\textbf{分析：}先根据已知条件和圆台结构特征分别求出两圆台的高，再根据圆台的体积公式直接代入计算即可得解。

\textbf{详解：}由题可得两个圆台的高分别为 $h_{\text{甲}}=\sqrt{[2(r_1-r_2)]^2-(r_1-r_2)^2}=\sqrt{3}(r_1-r_2)$，$h_{\text{乙}}=\sqrt{[3(r_1-r_2)]^2-(r_1-r_2)^2}=2\sqrt{2}(r_1-r_2)$，所以 $\dfrac{V_{\text{甲}}}{V_{\text{乙}}}=\frac{\frac{1}{3}(S_2+S_1+\sqrt{S_2S_1})h_{\text{甲}}}{\frac{1}{3}(S_2+S_1+\sqrt{S_2S_1})h_{\text{乙}}}=\frac{h_{\text{甲}}}{h_{\text{乙}}}=\frac{\sqrt{3}(r_1-r_2)}{2\sqrt{2}(r_1-r_2)}=\frac{\sqrt{6}}{4}$。
\end{solution}
\end{question}

\begin{question}[points=5]
已知 $a>1$，$\dfrac{1}{\log_8 a}-\dfrac{1}{\log_a 4}=-\dfrac{5}{2}$，则 $a=$ 。
\fillin[{$64$}]
\begin{solution}
\textbf{答案：}$64$

\textbf{分析：}将 $\log_8 a,\log_a 4$ 利用换底公式转化成 $\log_2 a$ 来表示即可求解。

\textbf{详解：}由题 $\frac{1}{\log_8 a}-\frac{1}{\log_a 4}=\frac{3}{\log_2 a}-\frac{1}{2}\log_2 a=-\frac{5}{2}$，整理得 $(\log_2 a)^2-5\log_2 a-6=0$，$\Rightarrow\log_2 a=-1$ 或 $\log_2 a=6$，又 $a>1$，所以 $\log_2 a=6=\log_2 2^6$，故 $a=2^6=64$。
\end{solution}
\end{question}

\begin{question}[points=5]
有6个相同的球，分别标有数字1、2、3、4、5、6，从中不放回地随机抽取3次，每次取1个球。记 $m$ 为前两次取出的球上数字的平均值，$n$ 为取出的三个球上数字的平均值，则 $|m-n|$ 不超过 $\dfrac{1}{2}$ 的概率是。
\fillin[{$\dfrac{7}{15}$}]
\begin{solution}
\textbf{答案：}$\dfrac{7}{15}$

\textbf{分析：}根据排列可求基本事件的总数，设前两个球的号码为 $a,b$，第三个球的号码为 $c$，则 $|a+b-3|\le 2c\le a+b+3$，就 $c$ 的不同取值分类讨论后可求随机事件的概率。

\textbf{详解：}从6个不同的球中不放回地抽取3次，共有 $A_6^3=120$ 种，设前两个球的号码为 $a,b$，第三个球的号码为 $c$，则 $\left|\frac{a+b+c}{3}-\frac{a+b}{2}\right|\le\frac{1}{2}$，故 $|2c-(a+b)|\le3$，故 $-3\le 2c-(a+b)\le3$，故 $a+b-3\le2c\le a+b+3$，若 $c=1$，则 $a+b\le5$，则 $(a,b)$ 为：$(2,3),(3,2)$，故有2种，若 $c=2$，则 $1\le a+b\le7$，则 $(a,b)$ 为：$(1,3),(1,4),(1,5),(1,6),(3,4),(3,1),(4,1),(5,1),(6,1),(4,3)$，故有10种，若 $c=3$，则 $3\le a+b\le9$，则 $(a,b)$ 为：$(1,2),(1,4),(1,5),(1,6),(2,4),(2,5),(2,6),(4,5),(2,1),(4,1),(5,1),(6,1),(4,2),(5,2),(6,2),(5,4)$，故有16种，若 $c=4$，则 $5\le a+b\le11$，同理有16种，若 $c=5$，则 $7\le a+b\le13$，同理有10种，若 $c=6$，则 $9\le a+b\le15$，同理有2种，共 $|m-n|$ 不超过 $\frac{1}{2}$ 时不同的抽取方法总数为 $2(2+10+16)=56$，故所求概率为 $\frac{56}{120}=\frac{7}{15}$。
\end{solution}
\end{question}

\section*{三、解答题}
共70分。解答应写出文字说明，证明过程或演算步骤。第17题~第21题为必考题，每个考题考生必须作答。第22、23题为选考题，考生根据要求作答。

\begin{question}[points=12]
某工厂进行生产线智能化升级改造，升级改造后，从该工厂甲、乙两个车间的产品中随机抽取150件进行检验，数据如下：

\begin{tabular}{c|ccc|c}

& 优级品 & 合格品 & 不合格品 & 总计 \\

\hline

甲车间 & 26 & 24 & 0 & 50 \\

乙车间 & 70 & 28 & 2 & 100 \\

\hline

总计 & 96 & 52 & 2 & 150 \\

\end{tabular}

（1）填写如下列联表：

\begin{tabular}{c|c|c}

& 优级品 & 非优级品 \\

\hline

甲车间 & & \\

乙车间 & & \\

\end{tabular}

能否有95%的把握认为甲、乙两车间产品的优级品率存在差异？能否有99%的把握认为甲，乙两车间产品的优级品率存在差异？

（2）已知升级改造前该工厂产品的优级品率 $p=0.5$，设 $\hat{p}$ 为升级改造后抽取的 $n$ 件产品的优级品率。如果 $\hat{p}>p+1.65\sqrt{\frac{p(1-p)}{n}}$，则认为该工厂产品的优级品率提高了，根据抽取的150件产品的数据，能否认为生产线智能化升级改造后，该工厂产品的优级品率提高了？（$\sqrt{150}\approx12.247$）

附：$K^2=\frac{n(ad-bc)^2}{(a+b)(c+d)(a+c)(b+d)}$，

\begin{tabular}{c|ccc}

$P(K^2\ge k)$ & 0.050 & 0.010 & 0.001 \\

\hline

$k$ & 3.841 & 6.635 & 10.828 \\

\end{tabular}
\begin{solution}
\textbf{答案：}（1）答案见详解；（2）答案见详解

\textbf{分析：}（1）根据题中数据完善列联表，计算 $K^2$，并与临界值对比分析；（2）用频率估计概率可得 $\hat{p}=0.64$，根据题意计算 $p+1.65\sqrt{\frac{p(1-p)}{n}}$，结合题意分析判断。

\textbf{详解：}（1）根据题意可得列联表：
\begin{tabular}{c|c|c}
& 优级品 & 非优级品 \\
\hline
甲车间 & 26 & 24 \\
乙车间 & 70 & 30 \\
\end{tabular}
可得 $K^2=\frac{150\times(26\times30-24\times70)^2}{50\times100\times96\times54}=\frac{75}{16}=4.6875$，因为 $3.841<4.6875<6.635$，所以有95%的把握认为甲、乙两车间产品的优级品率存在差异，没有99%的把握认为甲，乙两车间产品的优级品率存在差异。
（2）由题意可知：生产线智能化升级改造后，该工厂产品的优级品的频率为 $\frac{96}{150}=0.64$，用频率估计概率可得 $\hat{p}=0.64$，又因为升级改造前该工厂产品的优级品率 $p=0.5$，则 $p+1.65\sqrt{\frac{p(1-p)}{n}}=0.5+1.65\sqrt{\frac{0.5\times0.5}{150}}\approx0.5+1.65\times\frac{0.5}{12.247}\approx0.568$，可知 $\hat{p}>p+1.65\sqrt{\frac{p(1-p)}{n}}$，所以可以认为生产线智能化升级改造后，该工厂产品的优级品率提高了。
\end{solution}
\end{question}

\begin{question}[points=12]
记 $S_n$ 为数列 $\{a_n\}$ 的前 $n$ 项和，且 $4S_n=3a_n+4$。

（1）求 $\{a_n\}$ 的通项公式；

（2）设 $b_n=(-1)^{n-1}na_n$，求数列 $\{b_n\}$ 的前 $n$ 项和 $T_n$。
\begin{solution}
\textbf{答案：}（1）$a_n=4\cdot(-3)^{n-1}$；（2）$T_n=(2n-1)\cdot3^n+1$

\textbf{分析：}（1）利用退位法可求 $\{a_n\}$ 的通项公式；（2）利用错位相减法可求 $T_n$。

\textbf{详解：}（1）当 $n=1$ 时，$4S_1=4a_1=3a_1+4$，解得 $a_1=4$。当 $n\ge2$ 时，$4S_{n-1}=3a_{n-1}+4$，所以 $4S_n-4S_{n-1}=4a_n=3a_n-3a_{n-1}$ 即 $a_n=-3a_{n-1}$，而 $a_1=4\neq0$，故 $a_n\neq0$，故 $\frac{a_n}{a_{n-1}}=-3$，∴数列 $\{a_n\}$ 是以4为首项，$-3$ 为公比的等比数列，所以 $a_n=4\cdot(-3)^{n-1}$。
（2）$b_n=(-1)^{n-1}\cdot n\cdot4\cdot(-3)^{n-1}=4n\cdot3^{n-1}$，所以 $T_n=b_1+b_2+b_3+\cdots+b_n=4\cdot3^0+8\cdot3^1+12\cdot3^2+\cdots+4n\cdot3^{n-1}$，故 $3T_n=4\cdot3^1+8\cdot3^2+12\cdot3^3+\cdots+4n\cdot3^n$，所以 $-2T_n=4+4\cdot3^1+4\cdot3^2+\cdots+4\cdot3^{n-1}-4n\cdot3^n=4+4\cdot\frac{3(1-3^{n-1})}{1-3}-4n\cdot3^n=4+2\cdot3\cdot(3^{n-1}-1)-4n\cdot3^n=(2-4n)\cdot3^n-2$，$\therefore T_n=(2n-1)\cdot3^n+1$。
\end{solution}
\end{question}

\begin{question}[points=12]
如图，在以 $A,B,C,D,E,F$ 为顶点的五面体中，四边形 $ABCD$ 与四边形 $ADEF$ 均为等腰梯形，$BC\parallel AD$，$EF\parallel AD$，$AD=4$，$AB=BC=EF=2$，$ED=\sqrt{10}$，$FB=2\sqrt{3}$，$M$ 为 $AD$ 的中点。

（1）证明：$BM\parallel$ 平面 $CDE$；

（2）求二面角 $F-BM-E$ 的正弦值。

\begin{center}
  \includegraphics[width=.38\linewidth]{question-19-1.jpg}
\end{center}
\begin{solution}
\textbf{答案：}（1）证明见详解；（2）$\dfrac{4\sqrt{3}}{13}$

\textbf{分析：}（1）结合已知易证四边形 $BCDM$ 为平行四边形，可证 $BM\parallel CD$，进而得证；（2）作 $BO\perp AD$ 交 $AD$ 于 $O$，连接 $OF$，易证 $OB,OD,OF$ 三垂直，采用建系法结合二面角夹角余弦公式即可求解。

\textbf{详解：}（1）因为 $BC\parallel AD$，$EF=2$，$AD=4$，$M$ 为 $AD$ 的中点，所以 $BC\parallel MD$，$BC=MD$，四边形 $BCDM$ 为平行四边形，所以 $BM\parallel CD$，又因为 $BM\not\subset$ 平面 $CDE$，$CD\subset$ 平面 $CDE$，所以 $BM\parallel$ 平面 $CDE$。
（2）如图所示，作 $BO\perp AD$ 交 $AD$ 于 $O$，连接 $OF$，因为四边形 $ABCD$ 为等腰梯形，$BC\parallel AD$，$AD=4$，$AB=BC=2$，所以 $CD=2$，结合（1）$BCDM$ 为平行四边形，可得 $BM=CD=2$，又 $AM=2$，所以 $\triangle ABM$ 为等边三角形，$O$ 为 $AM$ 中点，所以 $OB=\sqrt{3}$，又因为四边形 $ADEF$ 为等腰梯形，$M$ 为 $AD$ 中点，所以 $EF=MD$，$EF\parallel MD$，四边形 $EFMD$ 为平行四边形，$FM=ED=AF$，所以 $\triangle AFM$ 为等腰三角形，$\triangle ABM$ 与 $\triangle AFM$ 底边上中点 $O$ 重合，$OF\perp AM$，$OF=\sqrt{AF^2-AO^2}=3$，因为 $OB^2+OF^2=BF^2$，所以 $OB\perp OF$，所以 $OB,OD,OF$ 互相垂直，以 $OB$ 方向为 $x$ 轴，$OD$ 方向为 $y$ 轴，$OF$ 方向为 $z$ 轴，建立 $O-xyz$ 空间直角坐标系，$F(0,0,3)$，$B(\sqrt{3},0,0)$，$M(0,1,0)$，$E(0,2,3)$，$\overrightarrow{BM}=(-\sqrt{3},1,0)$，$\overrightarrow{BF}=(-\sqrt{3},0,3)$，$\overrightarrow{BE}=(-\sqrt{3},2,3)$，设平面 $BFM$ 的法向量为 $\vec{m}=(x_1,y_1,z_1)$，平面 $EMB$ 的法向量为 $\vec{n}=(x_2,y_2,z_2)$，则 $\begin{cases} \vec{m}\cdot\overrightarrow{BM}=0 \\ \vec{m}\cdot\overrightarrow{BF}=0 \end{cases}$，即 $\begin{cases} -\sqrt{3}x_1+y_1=0 \\ -\sqrt{3}x_1+3z_1=0 \end{cases}$，令 $x_1=\sqrt{3}$，得 $y_1=3,z_1=1$，即 $\vec{m}=(\sqrt{3},3,1)$，$\begin{cases} \vec{n}\cdot\overrightarrow{BM}=0 \\ \vec{n}\cdot\overrightarrow{BE}=0 \end{cases}$，即 $\begin{cases} -\sqrt{3}x_2+y_2=0 \\ -\sqrt{3}x_2+2y_2+3z_2=0 \end{cases}$，令 $x_2=\sqrt{3}$，得 $y_2=3,z_2=-1$，即 $\vec{n}=(\sqrt{3},3,-1)$，$\cos\langle\vec{m},\vec{n}\rangle=\frac{\vec{m}\cdot\vec{n}}{|\vec{m}||\vec{n}|}=\frac{11}{\sqrt{13}\cdot\sqrt{13}}=\frac{11}{13}$，则 $\sin\langle\vec{m},\vec{n}\rangle=\frac{4\sqrt{3}}{13}$，故二面角 $F-BM-E$ 的正弦值为 $\frac{4\sqrt{3}}{13}$。
\end{solution}
\end{question}

\begin{question}[points=12]
设椭圆 $C:\frac{x^2}{a^2}+\frac{y^2}{b^2}=1(a>b>0)$ 的右焦点为 $F$，点 $M\left(1,\frac{3}{2}\right)$ 在 $C$ 上，且 $MF\perp x$ 轴。

（1）求 $C$ 的方程；

（2）过点 $P(4,0)$ 的直线与 $C$ 交于 $A,B$ 两点，$N$ 为线段 $FP$ 的中点，直线 $NB$ 交直线 $MF$ 于点 $Q$，证明：$AQ\perp y$ 轴。
\begin{solution}
\textbf{答案：}（1）$\dfrac{x^2}{4}+\dfrac{y^2}{3}=1$；（2）证明见解析

\textbf{分析：}（1）设 $F(c,0)$，根据 $M$ 的坐标及 $MF\perp x$ 轴可求基本量，故可求椭圆方程；（2）设 $AB:y=k(x-4)$，$A(x_1,y_1)$，$B(x_2,y_2)$，联立直线方程和椭圆方程，用 $A,B$ 的坐标表示 $y_1-y_Q$，结合韦达定理化简前者可得 $y_1-y_Q=0$，故可证 $AQ\perp y$ 轴。

\textbf{详解：}（1）设 $F(c,0)$，由题设有 $c=1$ 且 $\frac{b^2}{a}=\frac{3}{2}$，故 $\frac{a^2-1}{a}=\frac{3}{2}$，故 $a=2$，故 $b=\sqrt{3}$，故椭圆方程为 $\frac{x^2}{4}+\frac{y^2}{3}=1$。
（2）直线 $AB$ 的斜率必定存在，设 $AB:y=k(x-4)$，$A(x_1,y_1)$，$B(x_2,y_2)$，由 $\begin{cases} 3x^2+4y^2=12 \\ y=k(x-4) \end{cases}$ 可得 $(3+4k^2)x^2-32k^2x+64k^2-12=0$，故 $\Delta=1024k^4-4(3+4k^2)(64k^2-12)>0$，故 $-\frac{1}{2}<k<\frac{1}{2}$，又 $x_1+x_2=\frac{32k^2}{3+4k^2}$，$x_1x_2=\frac{64k^2-12}{3+4k^2}$，而 $N\left(\frac{5}{2},0\right)$，故直线 $BN:y=\frac{y_2}{x_2-\frac{5}{2}}\left(x-\frac{5}{2}\right)$，故 $y_Q=\frac{-3y_2}{2x_2-5}$，所以 $y_1-y_Q=y_1+\frac{3y_2}{2x_2-5}=\frac{y_1\times(2x_2-5)+3y_2}{2x_2-5}=\frac{k(x_1-4)\times(2x_2-5)+3k(x_2-4)}{2x_2-5}=k\frac{2x_1x_2-5(x_1+x_2)+8}{2x_2-5}=k\frac{2\times\frac{64k^2-12}{3+4k^2}-5\times\frac{32k^2}{3+4k^2}+8}{2x_2-5}=k\frac{128k^2-24-160k^2+24+32k^2}{3+4k^2}\cdot\frac{1}{2x_2-5}=0$，故 $y_1=y_Q$，即 $AQ\perp y$ 轴。
\end{solution}
\end{question}

\begin{question}[points=12]
已知函数 $f(x)=(1-ax)\ln(1+x)-x$。

（1）当 $a=-2$ 时，求 $f(x)$ 的极值；

（2）当 $x\ge0$ 时，$f(x)\ge0$ 恒成立，求 $a$ 的取值范围。
\begin{solution}
\textbf{答案：}（1）极小值为0，无极大值；（2）$a\le -\dfrac{1}{2}$

\textbf{分析：}（1）求出函数的导数，根据导数的单调性和零点可求函数的极值；（2）求出函数的二阶导数，就 $a\le -\frac{1}{2}$、$-\frac{1}{2}<a<0$、$a\ge0$ 分类讨论后可得参数的取值范围。

\textbf{详解：}（1）当 $a=-2$ 时，$f(x)=(1+2x)\ln(1+x)-x$，故 $f'(x)=2\ln(1+x)+\frac{1+2x}{1+x}-1=2\ln(1+x)-\frac{1}{1+x}+1$，因为 $y=2\ln(1+x),y=-\frac{1}{1+x}+1$ 在 $(-1,+\infty)$ 上为增函数，故 $f'(x)$ 在 $(-1,+\infty)$ 上为增函数，而 $f'(0)=0$，故当 $-1<x<0$ 时，$f'(x)<0$，当 $x>0$ 时，$f'(x)>0$，故 $f(x)$ 在 $x=0$ 处取极小值且极小值为 $f(0)=0$，无极大值。
（2）$f'(x)=-a\ln(1+x)+\frac{1-ax}{1+x}-1=-a\ln(1+x)-\frac{(a+1)x}{1+x},x>0$，设 $s(x)=-a\ln(1+x)-\frac{(a+1)x}{1+x},x>0$，则 $s'(x)=-\frac{a}{x+1}-\frac{a+1}{(1+x)^2}=-\frac{ax+2a+1}{(1+x)^2}$，当 $a\le -\frac{1}{2}$ 时，$s'(x)>0$，故 $s(x)$ 在 $(0,+\infty)$ 上为增函数，故 $s(x)>s(0)=0$，即 $f'(x)>0$，所以 $f(x)$ 在 $[0,+\infty)$ 上为增函数，故 $f(x)\ge f(0)=0$。当 $-\frac{1}{2}<a<0$ 时，当 $0<x<-\frac{2a+1}{a}$ 时，$s'(x)<0$，故 $s(x)$ 在 $(0,-\frac{2a+1}{a})$ 上为减函数，故在 $(0,-\frac{2a+1}{a})$ 上 $s(x)<s(0)$，即在 $(0,-\frac{2a+1}{a})$ 上 $f'(x)<0$ 即 $f(x)$ 为减函数，故在 $(0,-\frac{2a+1}{a})$ 上 $f(x)<f(0)=0$，不合题意，舍。当 $a\ge0$，此时 $s'(x)<0$ 在 $(0,+\infty)$ 上恒成立，同理可得在 $(0,+\infty)$ 上 $f(x)<f(0)=0$ 恒成立，不合题意，舍；综上，$a\le -\frac{1}{2}$。
\end{solution}
\end{question}

\begin{question}[points=10]
[选修4-4：坐标系与参数方程]

在平面直角坐标系 $xOy$ 中，以坐标原点 $O$ 为极点，$x$ 轴的正半轴为极轴建立极坐标系，曲线 $C$ 的极坐标方程为 $\rho=\rho\cos\theta+1$。

（1）写出 $C$ 的直角坐标方程；

（2）设直线 $l:\begin{cases} x=t \\ y=t+a \end{cases}$（$t$ 为参数），若 $C$ 与 $l$ 相交于 $A,B$ 两点，若 $|AB|=2$，求 $a$ 的值。
\begin{solution}
\textbf{答案：}（1）$y^2=2x+1$；（2）$a=\dfrac{3}{4}$

\textbf{分析：}（1）根据 $\begin{cases} \rho=\sqrt{x^2+y^2} \\ \rho\cos\theta=x \end{cases}$ 可得 $C$ 的直角方程；（2）将直线的新的参数方程代入 $C$ 的直角方程，法1：结合参数 $s$ 的几何意义可得关于 $a$ 的方程，从而可求参数 $a$ 的值；法2：将直线的直角方程与曲线的直角方程联立，结合弦长公式可求 $a$ 的值。

\textbf{详解：}（1）由 $\rho=\rho\cos\theta+1$，将 $\begin{cases} \rho=\sqrt{x^2+y^2} \\ \rho\cos\theta=x \end{cases}$ 代入，故可得 $\sqrt{x^2+y^2}=x+1$，两边平方后可得曲线的直角坐标方程为 $y^2=2x+1$。
（2）对于直线 $l$ 的参数方程消去参数 $t$，得直线的普通方程为 $y=x+a$。法1：直线 $l$ 的斜率为1，故倾斜角为 $\frac{\pi}{4}$，故直线的参数方程可设为 $\begin{cases} x=\frac{\sqrt{2}}{2}s \\ y=a+\frac{\sqrt{2}}{2}s \end{cases}$，$s\in\mathbb{R}$。将其代入 $y^2=2x+1$ 中得 $s^2+2\sqrt{2}(a-1)s+2(a^2-1)=0$，设 $A,B$ 两点对应的参数分别为 $s_1,s_2$，则 $s_1+s_2=-2\sqrt{2}(a-1)$，$s_1s_2=2(a^2-1)$，且 $\Delta=8(a-1)^2-8(a^2-1)=16-16a>0$，故 $a<1$，$\therefore |AB|=|s_1-s_2|=\sqrt{(s_1+s_2)^2-4s_1s_2}=\sqrt{8(a-1)^2-8(a^2-1)}=2$，解得 $a=\frac{3}{4}$。法2：联立 $\begin{cases} y=x+a \\ y^2=2x+1 \end{cases}$，得 $x^2+(2a-2)x+a^2-1=0$，$\Delta=(2a-2)^2-4(a^2-1)=-8a+8>0$，解得 $a<1$，设 $A(x_1,y_1),B(x_2,y_2)$，$\therefore x_1+x_2=2-2a$，$x_1x_2=a^2-1$，则 $|AB|=\sqrt{1+1^2}\cdot\sqrt{(x_1+x_2)^2-4x_1x_2}=\sqrt{2}\cdot\sqrt{(2-2a)^2-4(a^2-1)}=2$，解得 $a=\frac{3}{4}$。
\end{solution}
\end{question}

\begin{question}[points=10]
[选修4-5：不等式选讲]

实数 $a,b$ 满足 $a+b\ge3$。

（1）证明：$2a^2+2b^2>a+b$；

（2）证明：$|a-2b^2|+|b-2a^2|\ge6$。
\begin{solution}
\textbf{答案：}（1）证明见解析；（2）证明见解析

\textbf{分析：}（1）直接利用 $2a^2+2b^2\ge(a+b)^2$ 即可证明；（2）根据绝对值不等式并结合（1）中结论即可证明。

\textbf{详解：}（1）因为 $2a^2+2b^2-(a+b)=a^2-2ab+b^2=(a-b)^2\ge0$，当 $a=b$ 时等号成立，则 $2a^2+2b^2\ge(a+b)^2$，因为 $a+b\ge3$，所以 $2a^2+2b^2\ge(a+b)^2>a+b$。
（2）$|a-2b^2|+|b-2a^2|\ge|(a-2b^2)+(b-2a^2)|=|2a^2+2b^2-(a+b)|=2a^2+2b^2-(a+b)\ge(a+b)^2-(a+b)=(a+b)(a+b-1)\ge3\times2=6$。
\end{solution}
\end{question}

